% ===== 原文 pp.61--66 =====
% [原文 p.61]
背景场的小规范变换为
\begin{equation}
  A\to A+\dd\Lambda\,, \tag{4.3}
\end{equation}
其中 $\Lambda$ 是时空上的 0-形式。我们还有大规范变换，它们使和乐 $\oint A$ 平移整数。

对于具有 $G^{(0)}$ 0-形式对称性的理论，关联函数是背景场 $A$ 的函数。在存在背景场 $A$ 时，配分函数记作 $Z[A]$。

背景场的存在使作用量增加如下耦合
\begin{equation}
  2\pi\int A\wedge j^{d-1}\,, \tag{4.4}
\end{equation}
其中 $j^{d-1}$ 是 $G^{(0)}=\U(1)$ 0-形式对称性的 Noether 流。

% [Tutor Note: p.61 式 (4.4) 的 $2\pi$ 前置因子与后文 (4.14) 同源：把 (4.4) 中的背景 $A$ 换成电 1-形式背景 $B^e_2$、Noether 流取 $j^{d-p-1}=\star F$（$p=1$ 时即 $j^{d-2}=\star F$），即得 (4.14) 的 $2\pi\int B^e_2\wedge\star F$，说明两处采用同一归一化；大规范变换下 $\Delta S=2\pi\int\dd\Lambda\wedge j^{d-1}$ 是 $2\pi$ 的整数倍（电荷量子化），故 $Z$ 不变。]

\para{连续的 $p$-形式对称性。} 由于 $p\ge1$-形式对称性都是阿贝尔的，连续的 $p\ge1$-形式对称群必须具有如下形式
\begin{equation}
  G^{(p)}=U(1)^n\,,\qquad n>0\,. \tag{4.5}
\end{equation}
因此不失一般性，我们可以把讨论限制在
\begin{equation}
  G^{(p)}=U(1)\,, \tag{4.6}
\end{equation}
的 $p$-形式对称性上。在此情形下，我们有

\begin{definitionbox}{定义 4.1：连续 $p$-形式对称性的背景场}
连续 $G^{(p)}=\U(1)$ $p$-形式对称性的背景场是时空上的一个 $(p+1)$-形式 $B_{p+1}$，其场强
\begin{equation}
  F_{p+2}=\dd B_{p+1}\,, \tag{4.7}
\end{equation}
是时空上的闭 $(p+2)$-形式。
\end{definitionbox}

背景场的小规范变换为
\begin{equation}
  B_{p+1}\to B_{p+1}+\dd\Lambda_p\,, \tag{4.8}
\end{equation}
其中 $\Lambda_p$ 是时空上的 $p$-形式。我们还有大规范变换，它们使 $(p+1)$ 维闭子流形上的和乐 $\oint B_{p+1}$ 平移整数。

% [原文 p.62]
在具有式 (4.6) 对称性的理论中，关联函数是背景场 $B_{p+1}$ 的函数。特别地，存在 $B_{p+1}$ 时的配分函数表示为
\begin{equation}
  Z[B_{p+1}]\,. \tag{4.9}
\end{equation}
背景场的存在使作用量增加如下耦合
\begin{equation}
  2\pi\int B_{p+1}\wedge j^{d-p-1}\,, \tag{4.10}
\end{equation}
其中 $j^{d-p-1}$ 是 $G^{(p)}=\U(1)$ $p$-形式对称性的 Noether 流。

\begin{examplebox}{例 4.1：Maxwell 理论}
我们回到例 2.2 中引入的 Maxwell 理论。正如那里所讨论的，该理论具有电的
\begin{equation}
  G^{(1)}=U(1)\,, \tag{4.11}
\end{equation}
1-形式对称性，以及磁的
\begin{equation}
  G^{(d-3)}=U(1)\,, \tag{4.12}
\end{equation}
$(d-3)$-形式对称性。这些对称性的背景场分别是
\begin{equation}
  B^e_2\,,\qquad B^m_{d-2}\,, \tag{4.13}
\end{equation}
由 Noether 流的知识可知，在存在 $B^e_2$ 时，作用量包含项
\begin{equation}
  2\pi\int B^e_2\wedge\star F\,, \tag{4.14}
\end{equation}
而在存在 $B^m_{d-2}$ 时，作用量包含项
\begin{equation}
  2\pi\int B^m_{d-2}\wedge F\,. \tag{4.15}
\end{equation}
\end{examplebox}

% [Tutor Note: p.62 的量纲/次数核对：$B^e_2\wedge\star F$ 与 $B^m_{d-2}\wedge F$ 都是 $d$-形式，因此 (4.14)(4.15) 都能在 $d$ 维时空上积分；这说明电背景耦合到 $\star F$（即电 Noether 流 $j^{d-1}=\star F$，与 (4.4) 的归一化一致），磁背景则直接耦合到 $F$。]

\para{离散 $p$-形式对称性。} 现在考虑离散 $p\ge0$-形式对称群 $G^{(p)}$ 的情形。我们需要微分形式的离散版本。正如第 2.3 节所讨论的，这由时空上的离散上链提供。

\begin{definitionbox}{定义 4.2：离散 $p$-形式对称性的背景场}
离散 $G^{(p)}$ $p$-形式对称性的背景场是时空上的一个 $G^{(p)}$ 值 $(p+1)$-上链 $B_{p+1}$。
\end{definitionbox}

% [原文 p.63]
如果其他对称性的背景场没有被打开，那么 $B_{p+1}$ 是一个上闭链（cocycle），即
\begin{equation}
  \delta B_{p+1}=0\,. \tag{4.16}
\end{equation}
然而，如果其他对称性的背景场被打开，那么 $B_{p+1}$ 未必是上闭链。当 $G^{(p)}$ 参与非平凡的群扩张或高阶群对称性时，就会出现这种情况。我们稍后讨论这些方面。

背景场的规范变换为
\begin{equation}
  B_{p+1}\to B_{p+1}+\delta\Lambda_p\,, \tag{4.17}
\end{equation}
其中 $\Lambda_p$ 是时空上的 $G^{(p)}$ 值 $p$-上链。

\begin{examplebox}{例 4.2：离散规范理论}
考虑第 2.3 节所讨论的那种形式的离散（高阶形式）规范理论。它具有电的 $G^{(p)}$ $p$-形式对称性和磁的 $\widehat{G}^{(p)}$ $(d-p-1)$-形式对称性，二者都与 Noether 流的离散版本相联系。利用它们，我们可以把与背景场 $B^e_{p+1}$ 和 $B^m_{d-p}$ 的耦合描述为作用量中的如下各项
\begin{equation}
  2\pi\int\left(B^e_{p+1}\underset{\eta}{\cup}b_{d-p-1}+a_p\underset{\eta}{\cup}B^m_{d-p}\right)\,, \tag{4.18}
\end{equation}
其中回顾 $a_p$ 与 $b_{d-p-1}$ 是动力学规范场。
\end{examplebox}

\para{非阿贝尔规范理论的电 1-形式对称性} 考虑具有连通规范群 $\cG$ 的非阿贝尔规范理论。正如上面所讨论的，它的电 1-形式对称性是规范群中心的一个子群
\begin{equation}
  G^{(1)}\subseteq Z(\cG)\,. \tag{4.19}
\end{equation}
在打开 1-形式对称性的背景 $B^e_2$ 时，$\cG$ 规范理论的路径积分对所有那些取在群
\begin{equation}
  \cG/G^{(1)}\,, \tag{4.20}
\end{equation}
的丛上的规范场求和，对这些丛而言，提升为 $\cG$ 丛的障碍是
\begin{equation}
  w_2=B^e_2\,. \tag{4.21}
\end{equation}

% [原文 p.64]
\begin{examplebox}{例 4.3：纯 $\SU(2)$ Yang--Mills 理论}
作为一个具体例子，考虑纯 $\SU(2)$ Yang--Mills 理论，它具有
\begin{equation}
  G^{(1)}=Z(SU(2))=\Z_2\,, \tag{4.22}
\end{equation}
1-形式对称性。在存在背景 $B^e_2$ 时，路径积分对所有
\begin{equation}
  SU(2)/\Z_2=SO(3)\,, \tag{4.23}
\end{equation}
规范丛求和，其 Stiefel--Whitney 类 $w_2$ 被固定为
\begin{equation}
  w_2=B^e_2\,. \tag{4.24}
\end{equation}
\end{examplebox}

\para{非阿贝尔规范理论的磁 $(d-3)$-形式对称性} 回顾具有连通规范群 $\cG$ 的非阿贝尔规范理论具有由下式给出的磁 $(d-3)$-形式对称性
\begin{equation}
  G^{(d-3)}=\widehat{\cZ}\,, \tag{4.25}
\end{equation}
其中 $\cZ$ 是单连通群 $G$ 的中心 $Z(G)$ 的一个子群，而规范群 $\cG$ 通过 $G$ 表述为
\begin{equation}
  \cG=G/\cZ\,. \tag{4.26}
\end{equation}
这样的规范理论具有一个 $\cZ$ 值的动力学场 $w_2$，即把 $\cG$ 丛提升为 $G$ 丛的障碍类。在存在磁对称性的 $\widehat{\cZ}$ 值背景场 $B^m_{d-2}$ 时，作用量获得如下耦合
\begin{equation}
  2\pi\int B^m_{d-2}\underset{\eta}{\cup}w_2\,, \tag{4.27}
\end{equation}
其中 cup 积（上积）用到了自然的双同态
\begin{equation}
  \eta:\ \widehat{\cZ}\times\cZ\to\R/\Z\,. \tag{4.28}
\end{equation}

\begin{examplebox}{例 4.4：纯 $\SO(3)$ Yang--Mills 理论}
作为一个具体例子，考虑纯 $\SO(3)$ Yang--Mills 理论，它具有
\begin{equation}
  G^{(d-3)}=\Z_2\,. \tag{4.29}
\end{equation}
在此情形下，背景场开启如下耦合
\begin{equation}
  2\pi\int B^m_{d-2}\underset{\eta}{\cup}w_2=\pi\int B^m_{d-2}\cup w_2\,, \tag{4.30}
\end{equation}
其中 $w_2$ 是第二 Stiefel--Whitney 类，而右端是两个 $\Z_2$ 值上链之间的标准 cup 积。
\end{examplebox}

% [Tutor Note: 例 4.4 的方框在原文中跨 p.64--65 断开，末句 “where $w_2$ is the second Stiefel--Whitney class...” 位于 p.65 页顶；按规范已合并为一个 examplebox，故 (4.30) 及其说明连写。]

% [Tutor Note: p.64 式 (4.30) 把一般的双同态 $\eta:\widehat{\cZ}\times\cZ\to\R/\Z$ 具体化到 $\Z_2$：此时 $\eta(a,b)=\frac{1}{2}ab$，于是 $2\pi\int B^m_{d-2}\cup_\eta w_2=\pi\int B^m_{d-2}\cup w_2$，系数由 $2\pi\cdot\frac12=\pi$ 给出，可据此核对。]

% [原文 p.65]
\begin{figure}[htbp]
\centering
\includegraphics[width=0.35\textwidth]{fig12.png}
\figtitle{12}{$d=2$ 中两个三角剖分的局部片段，分别以红色和黑色画出，二者互为对偶。注意，其中一个三角剖分的 2-单形（面）对应于另一个三角剖分的 0-单形（顶点），意思是：一个三角剖分的每个面都恰好包含另一个三角剖分的一个顶点。类似地，1-单形（边）之间是一一对应的，意思是：一个三角剖分的一条边恰好与另一个三角剖分的一条边相交。}
\end{figure}

\para{背景场作为拓扑缺陷网络。} 对于离散 $G^{(p)}$（$p\ge0$）形式对称性，背景场 $B_{p+1}$ 可以用构成该 $G^{(p)}$ 对称性的、余维 $(p+1)$ 的拓扑缺陷网络来描述。这用到了 Poincaré 对偶的一种版本。

\begin{notebox}{算符 vs. 缺陷}
在讲义中，我们将等同地使用\textbf{算符}（operator）与\textbf{缺陷}（defect）这两个术语。这在相对论性 QFT 文献中是惯例。

其原因如下。缺陷可以理解为向理论中引入额外的自由度，这些自由度局域在空间的某个子流形上，从而改变 Hilbert 空间。现在考虑一个沿时空的子流形 $M_p$ 插入的 $p$ 维算符 $\cO$。我们可以选取一个时间方向，使得 $M_p$ 不完全位于某个时间切片之内。在这种情形下，$M_p$ 与时间切片的交叠可以看成一个缺陷，因为它在那个时间切片上一般地会改变 Hilbert 空间。
\end{notebox}

正如第 2.3 节所讨论的，上链 $B_{p+1}$ 用时空的一个三角剖分 $\cT$ 来定义。拓扑缺陷生活在一个对偶三角剖分 $\widehat{\cT}$ 上，后者定义如下：

% [原文 p.66]
\begin{itemize}[leftmargin=1.4em,itemsep=1pt,topsep=1pt]
\item $\widehat{\cT}$ 的一个顶点（或 0-单形）对应于 $\cT$ 的一个 $d$-单形，并置于该 $d$-单形内部。
\item $\widehat{\cT}$ 的一条边（或 1-单形）$\widehat{s}_1$ 对应于 $\cT$ 的一个 $(d-1)$-单形 $s_{d-1}$，并连接 $\widehat{\cT}$ 中与 $\cT$ 的两个 $d$-单形相对应的两个顶点，而这两个 $d$-单形共享该 $(d-1)$-单形 $s_{d-1}$。边 $\widehat{s}_1$ 与 $(d-1)$-单形 $s_{d-1}$ 恰好相交一次。
\item 依上述方式继续，$\widehat{\cT}$ 的一个 $i$-单形 $\widehat{s}_i$ 对应于 $\cT$ 的一个 $(d-i)$-单形 $s_{d-i}$，并与它相交于一点。$s_i$ 的边界由 $\widehat{\cT}$ 的 $(i-1)$-单形构成，而这些 $(i-1)$-单形对应于 $\cT$ 的 $(d-i+1)$-单形，后者的边界包含 $s_{d-i}$。
\end{itemize}

% [Tutor Note: p.66 第三条项目符号原文写 “The boundary of $s_i$”，未加 hat。按上下文，$\widehat{s}_i$ 是 $\widehat{\cT}$ 的 $i$-单形，其边界才由 $\widehat{\cT}$ 的 $(i-1)$-单形组成（对偶链的边界关系），因此这里应为 $\widehat{s}_i$。已用 700 dpi 裁图复核，原文确实印作 $s_i$，译文照录。]

关于 $d=2$ 中这两个三角剖分的图示，见图 12。

现在，拓扑缺陷网络按下述方式构造。对于 $\cT$ 的一个 $(p+1)$-单形 $s_{p+1}$，我们有
\begin{equation}
  B_{p+1}(s_{p+1})=g\in G^{(p)}\,. \tag{4.31}
\end{equation}
给定此式后，沿 $\widehat{\cT}$ 的对偶 $(d-p-1)$-单形 $\widehat{s}_{d-p-1}$ 插入余维 $(p+1)$ 的拓扑缺陷 $U_g$。

在存在背景场时，配分函数 $Z[B_{p+1}]$ 于是与相应拓扑缺陷网络的关联函数相同。在我们拥有离散 Noether 流 $j_{d-p-1}$（情况并非总是如此）的情形下，这很容易证明：因为在这种情形下，构成 $G^{(p)}$ $p$-形式对称性的拓扑余维 $(p+1)$ 算符可以表示为
\begin{equation}
  U_g=\int g(j_{d-p-1})\,,\qquad g\in G^{(p)}\,, \tag{4.32}
\end{equation}
其中 $j_{d-p-1}$ 是一个动力学的 $\widehat{G}^{(p)}$ 值 $(d-p-1)$-上链。如上所讨论，在这种情形下，打开背景场 $B_{p+1}$ 对应于向作用量加入耦合
\begin{equation}
  2\pi\int B_{p+1}\underset{\eta}{\cup}j_{d-p-1}\,. \tag{4.33}
\end{equation}
由 Poincaré 对偶，在此项存在时计算配分函数，恰好等价于计算置于与上链 $B_{p+1}$ Poincaré 对偶的余维 $(p+1)$ 链上的 $U_g$ 缺陷网络的关联函数。

% [Tutor Note: p.66 式 (4.32) 中 $U_g=\int g(j_{d-p-1})$ 的意思是：$g\in G^{(p)}$ 作为 $\widehat{G}^{(p)}$ 的特征（群同态 $\widehat{G}^{(p)}\to \U(1)$）作用在 $j$ 上再求和（指数化）。作者说 “This is easy to demonstrate” 所省略的论证正是 Poincaré 对偶：上链 $g(j)$ 与对偶链上的缺陷网络配对给出 (4.33)，可用 $d=2$、$G^{(0)}=\Z_2$ 的 Ising 模型作具体核对。]

\para{作为拓扑缺陷重排的规范变换。} 正如背景场描述时空上的拓扑缺陷网络，背景场的规范变换描述拓扑缺陷的形变，其中包括拓扑缺陷网络的局域重排。例子将在下一小节讨论，见图 13。
